
\documentclass{article} % For LaTeX2e
\usepackage{iclr2026_conference,times}

% Optional math commands from https://github.com/goodfeli/dlbook_notation.
\input{math_commands.tex}

\usepackage{hyperref}
\usepackage{url}


\title{Beyond Euclidean Spaces: Enhancing Retrieval-Augmented Generation through Manifold Representations}

% Authors must not appear in the submitted version. They should be hidden
% as long as the \iclrfinalcopy macro remains commented out below.
% Non-anonymous submissions will be rejected without review.

\author{Antiquus S.~Hippocampus, Natalia Cerebro \& Amelie P. Amygdale \thanks{ Use footnote for providing further information
about author (webpage, alternative address)---\emph{not} for acknowledging
funding agencies.  Funding acknowledgements go at the end of the paper.} \\
Department of Computer Science\\
Cranberry-Lemon University\\
Pittsburgh, PA 15213, USA \\
\texttt{\{hippo,brain,jen\}@cs.cranberry-lemon.edu} \\
\And
Ji Q. Ren \& Yevgeny LeNet \\
Department of Computational Neuroscience \\
University of the Witwatersrand \\
Joburg, South Africa \\
\texttt{\{robot,net\}@wits.ac.za} \\
\AND
Coauthor \\
Affiliation \\
Address \\
\texttt{email}
}

% The \author macro works with any number of authors. There are two commands
% used to separate the names and addresses of multiple authors: \And and \AND.
%
% Using \And between authors leaves it to \LaTeX{} to determine where to break
% the lines. Using \AND forces a linebreak at that point. So, if \LaTeX{}
% puts 3 of 4 authors names on the first line, and the last on the second
% line, try using \AND instead of \And before the third author name.

\newcommand{\fix}{\marginpar{FIX}}
\newcommand{\new}{\marginpar{NEW}}

%\iclrfinalcopy % Uncomment for camera-ready version, but NOT for submission.
\begin{document}


\maketitle

\begin{abstract}
Retrieval-Augmented Generation (RAG) fundamentally relies on 
vector similarity search to ground Large Language Models. However, 
conventional dense retrieval systems map textual data into 
Euclidean spaces, which severely limits their capacity to capture 
complex, hierarchical, or graph-like relationships inherent in 
real-world knowledge domains (e.g., medical ontologies or legal corpora). 
To address this limitation, we propose Manifold-RAG, a novel retrieval 
framework that shifts the representation space from standard Euclidean 
to non-Euclidean manifolds. Specifically, we project document chunks and 
user queries into a hyperbolic space (e.g., the Poincaré ball model), 
which mathematically accommodates tree-like topologies with minimal 
distortion. In our approach, we replace the standard cosine 
similarity with hyperbolic distance metrics during the dense retrieval 
phase. We anticipate that Manifold-RAG will significantly outperform 
traditional Euclidean baselines (such as BM25 and standard Sentence-BERT) 
on datasets exhibiting strong hierarchical structures, yielding 
substantial improvements in Recall@K and Mean Reciprocal Rank (MRR). 
Furthermore, we expect to demonstrate that this structural 
awareness achieves high retrieval accuracy even at lower embedding 
dimensions, thereby maintaining computational efficiency. By aligning 
the geometric properties of the embedding space with the underlying 
topology of the knowledge base, this work paves the way for more 
accurate and structurally-aware RAG systems, ultimately reducing 
hallucinations in complex knowledge domains.
\end{abstract}

\end{document}
